\pagebreak
\section{Pointer Initialization for Devices}
\label{sec:pointer_mapping}
\index{mapping!pointer}
\index{mapping!pointer attachment}
\index{pointer attachment}

Pointers that contain host addresses require that those addresses are
translated to device addresses for them to be useful in the context of a device
data environment. Broadly speaking, there are two scenarios where this is
important.

The first scenario is where the pointer is mapped to the device data
environment, such that references to the pointer inside a \kcode{target} region
refer to the corresponding pointer. \emph{Pointer attachment} ensures that the
corresponding pointer will contain a device address when all of the following
conditions are true:

\begin{itemize}
 \item the pointer is mapped by directive $A$ to a device;
 \item a list item that uses the pointer as its base pointer (call it the
       \emph{pointee}) is mapped, to the same device, by directive $B$, which may
       be the same as $A$; and
 \item the effect of directive $B$ is to create either the corresponding
       pointer or pointee in the device data environment of the device.
\end{itemize}

Given the above conditions, pointer attachment is initiated as a result of
directive $B$ and subsequent references to the pointee list item in a
\kcode{target} region that use the pointer will access the corresponding
pointee. The corresponding pointer remains in this \plc{attached} state until
it is removed from the device data environment.

The second scenario, which is only applicable for C/C++, is where:

\begin{itemize}
 \item the pointer is implicitly firstprivate inside a \kcode{target} construct
       when it appears as the base pointer to a list item on the construct; and
 \item the pointer does not appear explicitly as a list item in a \kcode{map}
       clause, \kcode{is_device_ptr} clause, or data-sharing attribute clause.
\end{itemize}

This scenario can be further split into two
cases: the list item is treated like an assumed-size array (e.g., \ucode{p[:]}
or \ucode{p[:0]}) or it is not.

If the list item is treated like an assumed-size array, this will trigger a
runtime check on entry to the \kcode{target} region for a previously mapped
list item (including list items mapped by the same construct) where the value
of the pointer falls within its \emph{mapped address range} or (if no such previously
mapped list item is found) its \emph{extended address range}. The mapped
address range corresponds to the actual mapped storage of a previously mapped
list item, while the extended address range includes the mapped address
range and additionally extends to the \emph{base address} of the previously
mapped list item (i.e., the first storage location of its base variable). For
instance, if the previously mapped list item is \ucode{x[5:n]}, the mapped
address range corresponds to the \ucode{n} elements starting from \ucode{x[5]},
while the extended address range includes all storage locations from
\ucode{x[0]} through the entire mapped address range. If such a match to a
previously mapped list item is found, the firstprivate pointer is initialized
to the device address corresponding to the value of the original pointer.
Otherwise, the firstprivate pointer retains its original value.

If the list item (again, call it the \emph{pointee}) is not treated like an
assumed-size array, the firstprivate pointer will be initialized such that
references in the \kcode{target} region to the pointee list item that use the
pointer will access the corresponding pointee (e.g., a reference to
\ucode{p[i]} will access the corresponding \ucode{p[i]} in the device data
environment).

The following example illustrates some basic concepts relating to pointer
initialization and attached pointers on a target device.

The pointers \ucode{ptr1} and \ucode{ptr2} each point to a dynamically
allocated array of \ucode{N} integers. The \kcode{target} construct maps these
arrays with the array section list items \ucode{ptr1[:N]} and \ucode{ptr2[:N]}.
As a comparison, note that the array \ucode{arr} can be mapped implicitly since
the compiler knows its length, while \ucode{ptr1[:N]} and \ucode{ptr2[:N]} must
be explicitly specified (at least when being mapped for the first time).

Additionally, the construct maps \ucode{ptr1}, while \ucode{ptr2} is treated as
firstprivate. Per the rules outlined above, \ucode{ptr1} becomes an attached
pointer that points to the corresponding \ucode{ptr1[:N]} array section, while
\ucode{ptr2} is initialized with the base address of the corresponding
\ucode{ptr2[:N]} array section. The OpenMP Specification imposes an additional restriction for
attached pointers like \ucode{ptr1} -- its value must not change between entry
and exit of the \kcode{target} region. The value of an attached pointer is
never updated back to the host as a result of a map or data-motion clause, so
this restriction prevents a situation in which the value of the pointer changes
on the device without a corresponding change occurring on the host. No such
restriction exists for firstprivate pointers that are initialized for the
device (like \ucode{ptr2}), since the semantics of firstprivate already ensures
that changes to the firstprivate copy will not be reflected in the original
copy.

The pointer \ucode{ptr3} is used inside the \kcode{target} construct, but it does
not appear in a data-mapping or data-sharing attribute clause. Nor is there a
\kcode{defaultmap} clause on the construct to indicate what its implicit
data-mapping or data-sharing attribute should be. For such a case, \ucode{ptr3}
will be implicitly treated as firstprivate within the construct and there will
be a runtime check to see if the host memory to which it is pointing has
corresponding memory in the device data environment, or is located within the
extended range of a mapped object. If this runtime check passes, the firstprivate
\ucode{ptr3} is initialized to point to the corresponding memory. But in this
case, the check does not pass and so it retains its original value (which is
overwritten by the result of the \ucode{malloc} call in the \kcode{target}
region). Since \ucode{ptr3} is firstprivate, the value to which it is assigned
in the \kcode{target} region is not returned into the original \ucode{ptr3} on
the host.

\cexample[5.0]{target_ptr_map}{1}

\index{directives!begin declare target@\kcode{begin declare target}}
\index{begin declare target directive@\kcode{begin declare target} directive}

In the following example, a \kcode{declare_target} directive applies to the global
pointer \ucode{p}.  Hence, the pointer \ucode{p} will exist on the device
for all \kcode{target} regions.

The pointer \ucode{p} is also used as the base pointer of an array section of a
\kcode{map} clause on a \kcode{target} construct. Since a corresponding pointer
\ucode{p} exists in the device data environment, on entry to the construct the
pointer becomes attached to the corresponding array section that results from the
map. On exit from the construct, no update is made to the host copy of \ucode{p}.
The implementation will continue to treat the device copy of \ucode{p} as if it
remains in an attached state, even though the pointee to which it was attached is
removed from the device. Consequentially, updates to or from the device copy of
\ucode{p} via a \kcode{target_update} directive or \kcode{map} clause will not
occur.

\cexample[5.1]{target_ptr_map}{2}

\index{directives!begin declare target@\kcode{begin declare target}}
\index{begin declare target directive@\kcode{begin declare target} directive}

The following two examples illustrate an incorrect and then correct reliance on
pointer attachment.

In example \example{target_ptr_map.3a}, the global pointer \ucode{p1} points to
array \ucode{x}, and the global pointer \ucode{p2} points to array \ucode{y} on
the host. The pointer \ucode{p1} appears in a \kcode{declare_target} directive.
The \kcode{target} construct will then implicitly map arrays \ucode{x} and
\ucode{y}, implicitly map \ucode{p1} (because of the \kcode{declare_target}
directive), and implicitly map (with the \kcode{storage} map type) the
assumed-size array \ucode{p2[:]}.

Inside the \kcode{target} construct, \ucode{p1} refers to the device copy
resulting from the \kcode{declare_target} directive and \ucode{p2} refers to a
predetermined firstprivate copy resulting from the implicit map of
\ucode{p2[:]}. However, \ucode{p1} still does not become an attached pointer to
\ucode{x} because the pointer attachment conditions listed at the beginning of
this section are not met. Specifically, there is no mapped list item for which
\ucode{p1} is a base pointer. Conversely, the firstprivate \ucode{p2} is
properly initialized to point to array \ucode{y} on entry to the \kcode{target}
construct. Note that prior to OpenMP 6.0 this initialization of \ucode{p2} was
not guaranteed, unless \ucode{y} was present on the device \emph{prior} the
encountering of the \kcode{target} construct. But OpenMP 6.0 expands the cases
where \ucode{p2} would be initialized to include this case where \ucode{y} is
created in the device data environment by the same construct. Consequentially,
the access to \ucode{p1[1]} inside the \kcode{target} construct will result in
undefined behavior, while the access to \ucode{p2[1]} will correctly access
\ucode{y[1]}.


\cexample[6.0]{target_ptr_map}{3a}

In example \example{target_ptr_map.3b}, a \kcode{target_data} construct is
added with an explicit map of \ucode{p1[:2]} to ensure pointer attachment for
the device copy of \ucode{p1}. Additionally, \ucode{y} is mapped on the
\kcode{target_data} construct to ensure that the firstprivate \ucode{p2}
created by the nested \kcode{target} construct is initialized to point to
\ucode{y} in implementations that support a pre-6.0 version of OpenMP.
Referencing \ucode{x} and \ucode{y} via either \ucode{p1} or \ucode{p2} inside
the \kcode{target} construct is therefore valid.

\cexample[6.0]{target_ptr_map}{3b}
%\clearpage

\index{routines!omp_target_is_accessible@\kcode{omp_target_is_accessible}}
\index{omp_target_is_accessible routine@\kcode{omp_target_is_accessible} routine}

In the following example, storage allocated on the host is not mapped in a
\kcode{target} region if it is determined that the host memory is accessible
from the device. On platforms that support host memory access from a target
device, it may be more efficient to omit \kcode{map} clauses and avoid the
potential memory allocation and data transfers that result from the map. The
\kcode{omp_target_is_accessible} API routine is used to determine if the host
storage of size \ucode{buf_size} is accessible on the device, and a
metadirective is used to select the appropriate directive variant (either a
\kcode{target} with a \kcode{map} clause or without one).

The \kcode{omp_target_is_accessible} routine will return true if the storage
indicated by the first and second arguments is accessible on the target device.
In this case, the host pointer \ucode{ptr} may be directly dereferenced in the
subsequent \kcode{target} region to access this storage, rather than mapping an
array section based off the pointer. Note that explicitly specifying the host
pointer in a \kcode{firstprivate} clause on the construct is not necessary as
of OpenMP 5.2. In OpenMP 5.1, removing the \kcode{firstprivate} clause will
result in an implicit presence check of the storage to which \ucode{ptr}
points, and since this storage is not mapped by the program, \ucode{ptr} will
be \bcode{NULL}-initialized in the \kcode{target} region. As of OpenMP 5.2, a
false presence check without the \kcode{firstprivate} clause will cause the
pointer to retain its original value.

\cexample[5.2]{target_ptr_map}{4}

\index{mapping!deep copy}
Similar to the previous example, the \kcode{omp_target_is_accessible} routine
is used to discover if a deep copy is required for the platform.  Here, the
\ucode{deep_copy} map, defined in the \kcode{declare mapper} directive, is used
if the host storage referenced by \ucode{s.ptr} (or \ucode{s\%ptr} in Fortran)
is not accessible from the device.

\cexample[5.2]{target_ptr_map}{5}
\ffreeexample[5.2]{target_ptr_map}{5}
